\section{Introduction}

The contributions of this chapter are heavily motivated by transaction systems.
Though each construction could be presented compactly and in isolation, the
contributions are clearer when explained in context.

We use the term transaction system to broadly capture all systems that involve
moving value. We describe the systems as if they handle just one type of
currency, but extensions to labelled assets are generally simple. Transaction
systems in this chapter must satisfy three properties. Balance means that no
money can be created by transactions. Theft prevention means that no user of the
system can spend money belonging to another. Privacy we define formally
in~\Cref{sec:txid-definitions}, and it is an intrinsic property of
the transaction systems we are working with.

The mechanism we give is agnostic to how the first two properties are achieved.
It constrains only the tag attached to a transaction, never its contents. So it
composes with any transaction system that already provides balance and theft
prevention. The confidential payload of a transaction in our scheme is just the
transaction of the underlying system, carried through untouched. What we add is
a way for a designated auditor to recover which of those payloads belong to a
given user, and only during a window the user has agreed to.

Confidentiality of wealth at rest is a trivial problem for cryptography, since
it can be solved by encrypting the value. Privacy of money in motion is the open
problem, and how much of it a system should offer has deep economic and
political roots that we do not attempt to settle here, deferring instead to the
relevant state, political, legal, and economic structures.
%(Domesday Book of 1086, Return of Owners of Land of 1873, and then Valuation
%Office Survey of 1910-1915, but significant portions of that are lost due to
%WWII bombings of the National Archives.)

The technical distinction that matters for us is not fiat versus non-fiat money,
but the treatment of identity. Fiat money is issued by an entity such as a
central bank, whose backing gives the currency its value and so centralises
issuance by design\footnote{Fiat is Latin for `let it be done'.}. Behaviour in
such systems is enforced not by protocol-level incentive compatibility but by
the ability of regulators and legal authorities to act against malicious
parties, which requires participants to be identifiable and so introduces the
requirement of registered identities. Cryptocurrencies instead replace the
issuer with a public ledger whose validators follow the protocol for economic
rather than institutional reasons. This is what lets them remain permissionless,
and lets them function with participants who have no registered identities at
all. They cover the entire spectrum of privacy, from pseudonymity with each
transaction in the clear to hidden values and anonymity sets equal to the total
number of users of the
system~\cite{bitcoin,ethereum,EPRINT:Noether15,SP:BCGGMT14}. Stable coins sit
between the two: cryptocurrencies with completely centralised issuance, backed
by fiat assets in an unofficial capacity and issued by private companies. By
trading volume and usage they, rather than central bank digital currencies
(CBDCs), are the largest digital assets backed by fiat~\cite{BIS26stablecoins}.

Our protocols generally deal with fiat money and so assume the role of a central
bank as issuer, though the issuer role is easy to distribute.

Cash is anonymous and tied to no identity at all, though depositing large sums
into the banking system often requires proof of provenance. The current digital
fiat system gives up this anonymity. Though transactions are not seen by the
central banks or other users, commercial banks see their clients' transactions
in the clear and tied to their identities. States can also compel commercial
banks to disclose this information in all its detail\footnote{United States
Congress, Bank Secrecy Act (1970) and Uniting and Strengthening America by
Providing Appropriate Tools Required to Intercept and Obstruct Terrorism (USA
PATRIOT) Act (2001)}.  Future digital fiat systems aim to improve the situation
with users proving ownership over the money being moved while the transactions
stay private. This is generally achieved through anonymous credentials.

\subsection{Related work}\label{sec:txid-related}

Prior work on regulatory oversight in privacy-preserving payment systems
typically supports either anonymity revocation or tracing.  Anonymity revocation
lets an auditor identify the user responsible for a specific transaction.
Tracing, by contrast, lets the auditor recover all transactions associated with
a given user. Tracing can be emulated by applying anonymity revocation to every
transaction in turn, but this costs time linear in the number of transactions
and exposes information about non-targeted users.

Anonymity revocation is achieved through attaching a verifiable encryption or
other hidden user identifier to each transaction, so that an auditor holding the
corresponding secret key can selectively recover the sender's
identity~\cite{SCN:CamHohLys06,CANS:BDET21,INDOCRYPT:CDLP22}.  Tracing a user's
activity this way requires scanning and decrypting every transaction on the
ledger, incurring overhead linear in the ledger size and exposing unrelated
users' activity to the auditor. These protocols enable identity recovery, but
not efficient targeted tracing.

Per-transaction tags can be derived from users' tracing keys combined with a
monotonically increasing counter, and use zero-knowledge proofs to show the tags
are correctly formed and bound to the transaction's
origin~\cite{CCS:KiaKohSar22,CANS:SarKiaKoh24}.  Given the tracing secret of a
target user, an auditor identifies the corresponding transactions by recomputing
the tags and matching them against the ledger, and the counter binding also
stops a user later repudiating their own transactions. These systems are
necessarily tied to users' identities.  All of these share the same limitation:
because the counter must increase by one with each transaction and leave no
gaps, a user's transactions must be generated sequentially, which prevents
concurrent transaction creation and degrades performance in high-throughput
settings.

Cryptocurrencies generally provide auditability through interactive protocols or
global consistency checks. ZKLedger lets auditors verify transaction
correctness, but requires interaction and incurs overhead linear in the number
of participants~\cite{EPRINT:NarVasVir18}.  Privacy-focused cryptocurrencies
such as Monero, Zcash, and Quisquis achieve strong anonymity guarantees but do
not directly support efficient
tracing~\cite{EPRINT:Noether15,SP:BCGGMT14,AC:FMMO19,FC:BAZB20}.  Retrofitting
tracing into such systems typically requires additional interaction, user
cooperation, or linear-time analysis~\cite{ACNS:PBGP25}.

Platypus and ZKAT incorporate regulatory policy enforcement, including
enforcement of balance limits, but neither provides a mechanism for efficient
per-user tracing across arbitrary
transactions~\cite{CCS:WKDC22,EPRINT:ACDDET19}.

\Cref{tab:txid-related} summarises the comparison. The column that matters is
the third: the cost of recovering one user's transactions for one epoch. For
revocation-based schemes this is linear in the size of the ledger $|L|$, because
every transaction must be decrypted before it is known whether it belongs to the
target. For our scheme it is $\Max$ pairings and $\Max$ lookups in a ledger
indexed by tag, where $\Max$ is the per-epoch transaction bound. At $\Max =
2^{10}$ this is under a second, and crucially it does not grow as the ledger
does, nor does it reveal anything about the transactions of any user other than
the target.

\begin{table}[t] \centering \caption{Comparison with prior approaches to
regulatory oversight. $|L|$ is the number of transactions on the ledger, $n$ the
number of registered users, and $\Max$ the per-epoch transaction bound.}
\label{tab:txid-related}
  %
  \small
  %
  \begin{tabular}{l l l c c} \toprule
    %
    Scheme & Mechanism & Trace one user & Concurrent & Expiring \\ & & &
submission & permission \\
    %
    \midrule
    %
    \cite{SCN:CamHohLys06,CANS:BDET21,INDOCRYPT:CDLP22} & revocation & $O(|L|)$
decryptions & yes & no \\
    %
    \cite{CCS:KiaKohSar22,CANS:SarKiaKoh24} & tags, sequential counter &
$O(|L|)$ tag matches & no & no \\
    %
    \cite{EPRINT:NarVasVir18} & interactive audit & $O(n)$, interactive & yes &
no \\
    %
    \cite{CCS:WKDC22,EPRINT:ACDDET19} & policy enforcement & not supported & yes
& n/a \\
    %
    \midrule
    %
    This chapter & tags, unordered index & $O(\Max)$ pairings & yes & yes \\
    %
    \bottomrule \end{tabular}
%
\end{table}

\paragraph{Our contribution.} We introduce a tracing mechanism that eliminates
the need for encrypted user identifiers and for the monotonically increasing
counters of prior tag-based schemes. In place of an ordered counter, a
transaction carries an index drawn from a bounded, epoch-local space $[0,\Max)$,
and the only constraint is that an index is not reused within an epoch. Order is
irrelevant, so a user may partition the space among concurrent workers and
submit transactions in parallel and in any order, which is what the sequential
counter of prior work forbids. Users keep fine-grained control through
epoch-bound tracing permissions. A permission lets an auditor identify every
transaction a specific user made within one window, and nothing outside it. It
expires with the epoch, and tracing the next one needs a fresh grant. We
formalise tracing as an ideal
functionality capturing anonymity, non-repudiation, and permission expiration,
give a black-box construction from signatures, commitments, pseudorandom
functions, one-way functions, and zero-knowledge proofs, and instantiate it
concretely in two ways. The first works in pairing-based groups with algebraic
pseudorandom functions and tailored sigma protocols. Its setup is transparent,
with no party sampling any group generator, and the BBS+ signature scheme lets
the registration authority certify a user's keys without ever learning their
tracing key. The second takes both pseudorandom functions to be hash functions
and proves the two corresponding relations with a general purpose zk-SNARK. That
buys constant-size transaction proofs and faster verification, and pays for them
with a trusted setup and a slower prover.

\paragraph{Layout.} \Cref{sec:txid-definitions} gives the syntax of a tracing
scheme and formalises its security as an ideal functionality.
\Cref{sec:txid-construction} gives a construction from generic building blocks.
\Cref{sec:txid-instantiation} gives the two concrete instantiations, proved
secure in \Cref{sec:txid-security}.  \Cref{sec:txid-evaluation} evaluates both,
and \Cref{sec:txid-summary} concludes.
